% SPDX-License-Identifier: Apache-2.0
\section{The Engines Against the Bridge They Really Talk Through}
\label{sec:x-dmawb}

\S\ref{sec:x-dma} and \S\ref{sec:x-dmaw} measure \code{LineFetch} and
\code{LineStore} against a memory the example defines: local, on the
peripheral side of the link, and answering a burst in full. That is
the right model for measuring an engine. It is not the path to the
board's memory.

That path goes through \code{AxiWb}, which puts each word of a burst
on a pipelined Wishbone as its own request. This run is the engines
behind a real one.

\subsection{The comment that was aimed at the wrong thing}

\code{ex\_dma} defines its own memory, and says of it:

\begin{quote}
A model that answered one beat per request would make a burst look
like a word and the run would prove nothing about bursting, which is
the whole of what is being measured here.
\end{quote}

That was written to justify making the \emph{test model} stricter.
The peripheral the engines were meant to be pointed at answered one
beat per request, which is the model that comment rejects as proving
nothing. The test had been hardened against a behaviour the system
was never checked for, and issue 471 is what the system turned out to
do.

So this example exists because an example is not the system. The
model was never wrong; it was answering a question nobody had asked
of the real path.

\subsection{What the run does}

The two engines are chained rather than checked apart. The fetch
engine reads a memory through one bridge, its words feed the store
engine, and that writes them to a second memory through a second
bridge. The check at the end, that the far memory holds what the near
one held, is therefore a check of both engines and both bridges at
once, and a word lost anywhere in the path is a word missing at the
end.

A bridge and a memory each rather than an arbiter between them: a
host sharing one link is the board's business, and one set of
Wishbone lines cannot carry two bridges.

Both memories stall for ten cycles before they answer anything and
answer after two thereafter, which is nearer a controller that
calibrates than a memory that answers at once.

\subsection{That the run ended is its own assertion}

The interesting failure here is not a wrong word. Both engines count
the beats of a burst themselves and neither reads the returned beat's
\code{last}, so a peripheral that answers too few beats leaves them
waiting for the rest for ever. The fault is a run that does not
finish.

Asserted separately for that reason. Without it the failure would be
reported as missing data and the reason would be a guess.

Against the bridge as it was before issue 471, with everything else
the same, the run reaches its cycle cap with both engines still
running and one word of sixty four landed. Against the bridge as it
is, sixty four words cross in 435 cycles and both engines go idle.

\subsection{What it costs, measured}

\S\ref{sec:x-dma} moves sixty four words in 93 cycles against a local
memory that answers at once. The same sixty four words through two
real bridges, chained through both engines, with a memory that
calibrates and then answers after two cycles, take 435.

Neither number is the other's competitor: the first measures an
engine and the second measures a path. The figure worth carrying is
that the path is where the cycles are, which is what a burst served
a word at a time costs, and what issuing several requests before
their acknowledgements would win back.

\lstinputlisting[style=ascii,caption={What \code{ex\_dmawb} printed
when the build ran it.},%
label={lst:x-dmawb-out}]{out_dmawb.txt}
